\chapter{乘积}\label{chap:products}

\noindent\textbf{第 \thechapter{} 章学习目标}
\begin{itemize}
  \item 证明仿射簇的乘积是仿射簇.
  \item 由 $\mathcal{I}(X)$ 和 $\mathcal{I}(Y)$ 确定消没理想 $\mathcal{I}(X\times Y)$.
  \item 描述仿射簇乘积的不可约性, 维数和光滑性与其各因子的关系.
  \item 描述 $R$-模张量积的元素, 并用张量积的规则判定两个元素何时相等.
  \item 在具体例子中, 将 $R$-模和 $R$-代数的张量积与更熟悉的模和代数认同.
  \item 将 $X\times Y$ 的坐标环计算为一个张量积.
\end{itemize}

从已有数学对象构造新对象的一种常见方法是取 Cartesian 积. 例如, 两个加法群的乘积配备逐分量加法后是一个群, 两个拓扑空间的乘积配备乘积拓扑后是一个拓扑空间. 本章试图理解两个仿射簇的 Cartesian 积.

要看出仿射簇的乘积本身也是仿射簇, 需要稍作论证. 做到这点后, 就可以研究乘积的几何性质, 例如不可约性, 维数和光滑性. 本章前两节专门发展这些内容, 最后的结论是: 两个光滑仿射簇 $X$ 和 $Y$ 的乘积是维数为 $\dim(X)+\dim(Y)$ 的光滑仿射簇.

熟悉了仿射簇乘积的几何性质后, 代数与几何的等价性自然引出一个问题: 与几何上的 Cartesian 积相对应的坐标环上的代数运算是什么? 换句话说, $K$-代数 $K[X\times Y]$ 与 $K[X]$ 和 $K[Y]$ 有何关系?

答案是 $K[X\times Y]$ 是 $K[X]$ 与 $K[Y]$ 的张量积. 在 $X=Y=\mathbb{A}^{1}$ 的特殊情形, 这就是说 $K[x,y]$ 是 $K[x]$ 与 $K[y]$ 的张量积. 对许多读者而言, 张量积可能是陌生的, 即便不陌生, 也可能令人望而生畏, 所以我们不假设读者事先了解它. 我们从一个目标出发: 定义一种代数运算, 将 $K[x]$ 和 $K[y]$ 合在一起得到 $K[x,y]$. 围绕这个目标逐步建立定义. 本章最后证明, 对任意仿射簇 $X$ 和 $Y$, 都有 $K[X\times Y]=K[X]\otimes_{K}K[Y]$.

\par\smallskip
{\footnotesize
\noindent\textcopyright{} The Author(s) 2025\par
\noindent E. Clader, D. Ross, \emph{Beginning in Algebraic Geometry}, Undergraduate Texts in Mathematics,
\href{https://doi.org/10.1007/978-3-031-88819-9_9}{\nolinkurl{https://doi.org/10.1007/978-3-031-88819-9_9}}.\par}
